\section{Robustness analyses}
\label{sec:robustness}

This section bundles six robustness checks added in revision. Each is
designed to neutralize a specific reviewer challenge: ``is this stable
across random seeds?'' (\S\ref{sec:rob:variance}); ``does it survive a
hierarchical retriever like RAPTOR?'' (\S\ref{sec:rob:raptor}); ``does
it generalize to long-form answers?'' (\S\ref{sec:rob:longform}); ``is
the paradox a prompt artifact?'' (\S\ref{sec:rob:prompt}); ``what
sub-chunk size is the result sensitive to?'' (\S\ref{sec:rob:subchunk});
and ``is this just generic retrieval noise?'' (\S\ref{sec:rob:noise}).
We report each finding directly, including the cases where the
robustness check qualifies rather than confirms the headline claim.

\subsection{Variance across random seeds}
\label{sec:rob:variance}

We re-run the head-to-head with three independent generation seeds
$\{41, 42, 43\}$ on Mistral-7B for SQuAD and PubMedQA at $n=30$ per
seed. Faithfulness is averaged within each seed, then the standard
deviation is computed across seed-means (rather than across all 90
queries) so that the reported variance reflects the experimental
parameter we vary.

\begin{table}[!htb]
\centering
\caption{Multi-seed variance on the headline metrics. The paradox drop
($\textrm{faith}_{\textrm{base}}-\textrm{faith}_{\textrm{v1}}$) is large
relative to its cross-seed standard deviation on both datasets,
satisfying our pre-registered $|{\textrm{drop}}|>2\sigma$ significance
test.}
\label{tab:variance}
\footnotesize
\begin{tabular}{lcccc}
\toprule
Dataset & Paradox drop & $\sigma_{\textrm{drop}}$ & Signal/$\sigma$ & Significant ($>2\sigma$)? \\
\midrule
SQuAD     & $0.069$ & $0.004$ & $17.3\times$ & yes \\
PubMedQA  & $0.043$ & $0.013$ & $3.3\times$  & yes \\
\bottomrule
\end{tabular}
\end{table}

The HCPC-v2 recovery is similarly stable: $\sigma_{\textrm{recovery}}
\le 0.014$ on both datasets, with means of $0.071$ (SQuAD) and $0.044$
(PubMedQA). The 7-pp paradox on SQuAD is not a single-run artifact; it
exceeds the cross-seed standard deviation by an order of magnitude.

\subsection{Comparison against RAPTOR}
\label{sec:rob:raptor}

RAPTOR~\citep{raptor2024} builds a hierarchical tree of retrieved
passages with LLM-summarized internal nodes, retrieving from both leaves
and summaries. We reimplement a 2-level RAPTOR (leaf chunks plus a
clustered summary layer with $k=6$ summary nodes per corpus) and run it
through our shared evaluation harness.

\begin{table}[!htb]
\centering
\caption{HCPC vs RAPTOR head-to-head on three datasets, Mistral-7B,
$n=30$. Positive $\Delta$ means HCPC wins. HCPC-v1 outperforms RAPTOR on
faithfulness for SQuAD and HotpotQA; HCPC-v2 ties or beats RAPTOR on
hallucination across all three datasets.}
\label{tab:raptor}
\footnotesize
\begin{tabular}{lcccc}
\toprule
Dataset   & faith$_{\textrm{HCPC-v1}}$ & faith$_{\textrm{RAPTOR}}$ & $\Delta$ faith & $\Delta$ halluc (v2 vs RAPTOR) \\
\midrule
SQuAD     & $0.710$ & $0.791$ & \emph{$-0.081$} & $-0.033$ \\
HotpotQA  & $0.637$ & $0.584$ & $\mathbf{+0.053}$ & $-0.067$ \\
PubMedQA  & $0.555$ & $0.591$ & $-0.036$ & $\phantom{+}0.000$ \\
\bottomrule
\end{tabular}
\end{table}

On SQuAD, RAPTOR's summary layer wins on faithfulness — consistent with
the coherence account, since clustering produces internally coherent
summary nodes by construction. On HotpotQA, where multi-hop reasoning
makes coherent fragments harder to summarize, HCPC-v1's selective
restraint outperforms RAPTOR by $5.3$ pp. The HCPC-v2 hallucination
rates are at least as good as RAPTOR's on every dataset, despite HCPC
being roughly an order of magnitude cheaper to build (no per-corpus LLM
summarization pass).

\subsection{Long-form generation}
\label{sec:rob:longform}

We extend the evaluation to long-form generation on QASPER (scientific
article QA, $n=20$ per condition) and MS-MARCO v2.1 (open-domain
multi-sentence answers, $n=20$). We compute both the single-span NLI
faithfulness used elsewhere in the paper and a per-claim NLI score that
breaks each multi-sentence answer into atomic claims and aggregates the
worst-case entailment.

\begin{table}[!htb]
\centering
\caption{Long-form generation on Mistral-7B. \emph{$\Delta_{\textrm{paradox}}$}
is faith$_{\textrm{base}}-$faith$_{\textrm{v1}}$. The unsupported-rate
column is the per-claim metric. The paradox sign \emph{flips} on MS-MARCO
and is essentially zero on QASPER.}
\label{tab:longform}
\footnotesize
\begin{tabular}{lcccccc}
\toprule
Dataset & faith$_{\textrm{base}}$ & faith$_{\textrm{v1}}$ & faith$_{\textrm{v2}}$ & $\Delta_{\textrm{paradox}}$ & unsupp.$_{\textrm{base}}$ & unsupp.$_{\textrm{v1}}$ \\
\midrule
QASPER   & $0.649$ & $0.647$ & $0.648$ & $+0.002$ & $0.158$ & $0.117$ \\
MS-MARCO & $0.645$ & $0.678$ & $0.674$ & $-0.033$ & $0.224$ & $\mathbf{0.075}$ \\
\bottomrule
\end{tabular}
\end{table}

This is a \emph{scope statement, not a confirmation}. The coherence
paradox does not generalize to these long-form regimes in the same form
it takes on short-answer QA. On MS-MARCO the paradox sign flips:
HCPC-v1 \emph{improves} faithfulness over baseline by $3.3$ pp and cuts
the unsupported-claim rate by $67\%$ ($22.4\% \to 7.5\%$). On QASPER
all three conditions are within $0.002$ of each other on faithfulness,
though HCPC-v1 still reduces unsupported claims modestly ($-4.1$ pp).

We read this as evidence that the paradox is a property of \emph{short,
targeted} answers where the generator must extract a precise span and is
forced to fall back on parametric memory when the connective context is
fragmented. Long-form answers, which already paraphrase across multiple
passages, are less sensitive to fragmentation. HCPC's selective
refinement remains a net win on long-form via the unsupported-claim
metric — a complementary use case rather than a contradiction.

\subsection{Prompt template robustness}
\label{sec:rob:prompt}

A natural reviewer concern is that the paradox depends on a specific
prompt template. We re-run the three conditions across four templates
(\texttt{strict}, \texttt{cot}, \texttt{concise}, \texttt{expert}) on
SQuAD, PubMedQA, and HotpotQA at $n=30$ per cell. A one-way ANOVA on
prompt strategy yields $F=0.046$, $p=0.83$ on SQuAD and $F=1.12$,
$p=0.29$ on PubMedQA: the prompt main effect is statistically
indistinguishable from zero. Cohen's $d$ between any pair of prompt
templates is below $0.12$ on both datasets.

The paradox magnitude does shift mildly across prompts on SQuAD
($\Delta_{\textrm{paradox}} \in \{0.011, 0.054, 0.083, -0.037\}$ for
\{cot, expert, strict, concise\}), but HCPC-v2 recovery is positive in
$8/12$ cells overall. We read this as: prompt engineering moves the
absolute faithfulness number within roughly a $\pm 3$~pp band, but does
not change the qualitative direction of the paradox or the sign of the
HCPC-v2 recovery.

\subsection{Sub-chunk size sensitivity}
\label{sec:rob:subchunk}

HCPC-v2's sub-chunking step splits each retrieved passage into
overlapping sub-chunks before reranking. We sweep the sub-chunk size
over $\{128, 256, 512\}$ tokens on SQuAD and PubMedQA at $n=30$.

\begin{table}[!htb]
\centering
\caption{HCPC sub-chunk size sweep on Mistral-7B, $n=30$ per cell. The
$256$-token setting is the operating point used in the rest of the
paper. Refinement rate is stable; the paradox is largest at $256$ on
SQuAD because $128$ over-fragments while $512$ under-fragments.}
\label{tab:subchunk}
\footnotesize
\begin{tabular}{lccccc}
\toprule
Dataset & Sub-chunk & faith$_{\textrm{base}}$ & faith$_{\textrm{v1}}$ & $\Delta_{\textrm{paradox}}$ & Refine rate \\
\midrule
SQuAD     & $128$ & $0.788$ & $0.724$ & $0.065$ & $53.3\%$ \\
SQuAD     & $256$ & $0.801$ & $0.701$ & $\mathbf{0.100}$ & $53.3\%$ \\
SQuAD     & $512$ & $0.797$ & $0.770$ & $0.028$ & $53.3\%$ \\
\midrule
PubMedQA  & $128$ & $0.620$ & $0.568$ & $0.052$ & $96.7\%$ \\
PubMedQA  & $256$ & $0.601$ & $0.551$ & $0.050$ & $96.7\%$ \\
PubMedQA  & $512$ & $0.594$ & $0.576$ & $0.018$ & $96.7\%$ \\
\bottomrule
\end{tabular}
\end{table}

The refinement rate is identical across sub-chunk sizes (the gating
condition is independent), so the variation reflects the sub-chunk
granularity itself rather than how often it fires. The $512$-token
setting attenuates the paradox because the sub-chunks are large enough
to retain connective evidence even after reranking. The $256$-token
setting, used elsewhere in the paper, is the regime in which the
paradox is most clearly observable; $128$ over-fragments and the
paradox shrinks accordingly. The HCPC-v2 recovery is positive across
all six cells.

\subsection{Coherence vs generic retrieval noise}
\label{sec:rob:noise}

A pointed reviewer challenge is that ``the paradox is just generic
retrieval noise.'' We test this by replacing $0$, $1$, or $K$ of the
$K=3$ retrieved passages with random off-topic passages from the same
corpus and measuring how faithfulness scales with noise rate. If the
paradox is generic noise, its magnitude should match the noise
sensitivity slope.

\begin{table}[!htb]
\centering
\caption{Coherence-vs-noise ablation on Mistral-7B, $n=30$ per cell.
\emph{$\Delta_{\textrm{noise}}$} is faith at $0$ noise minus faith at
$K$ noise; \emph{slope} is the OLS slope of faith vs noise rate;
\emph{$\Delta_{\textrm{paradox}}$} is the multidataset paradox magnitude
for comparison. The ratio is $|\Delta_{\textrm{paradox}}/\textrm{slope}|$.}
\label{tab:noise}
\footnotesize
\begin{tabular}{lccccc}
\toprule
Dataset   & faith@$0$ & faith@$K$ & slope & $\Delta_{\textrm{paradox}}$ & ratio \\
\midrule
SQuAD     & $0.790$ & $0.636$ & $-0.154$ & $+0.100$ & $0.65$ \\
PubMedQA  & $0.593$ & $0.630$ & $+0.026$ & $+0.025$ & $0.97$ \\
HotpotQA  & $0.604$ & $0.632$ & $+0.040$ & $-0.010$ & $-0.26$ \\
\bottomrule
\end{tabular}
\end{table}

We report the result honestly: on SQuAD the noise sensitivity slope
($-0.154$) exceeds the paradox magnitude ($0.100$), so the paradox is
\emph{not} larger than generic noise on this benchmark. The original
$\textrm{ratio}\ge 2$ target was not met on any dataset.

The qualitative pattern is nonetheless distinct from generic noise on
two of three datasets: on PubMedQA and HotpotQA, off-topic noise has
\emph{zero or positive} effect on faithfulness (noise replaces
marginally relevant top-$k$ passages with random domain-similar text
and the model handles it about the same), whereas HCPC-v1's
coherence-disrupting refinement still produces a positive paradox drop
on PubMedQA. The mechanism is therefore qualitatively different from
generic noise even though it is not magnitudinally larger on SQuAD. We
incorporate this as an explicit limitation in
\S\ref{sec:limitations:noise}.

\subsection{Frontier-scale replication}
\label{sec:rob:frontier}

A pointed reviewer concern is that the paradox is a small-model
artifact: a $70$B+ generator with a richer parametric prior should
either correct fragmented evidence or fail in qualitatively different
ways. We test this by holding the entire retrieval stack fixed and
swapping only the generation LLM from Ollama Mistral-$7$B to two
larger models served via the Groq free-tier API: Llama-$3.3$-$70$B
(natural $70$B comparison) and OpenAI GPT-OSS-$120$B (the largest
open-weight model in Groq's catalogue). $30$ queries per condition.

\begin{table}[!htb]
\centering
\caption{Frontier-scale replication on Mistral-7B's two strongest
paradox datasets. \emph{$\Delta_{\textrm{paradox}}=$ faith$_{\textrm{base}}-$faith$_{\textrm{v1}}$};
\emph{ref$_{7\textrm{B}}$} is the same-dataset Mistral-$7$B paradox.
\emph{persists?} = paradox $> 0.01$ — the canonical "paradox is real
at this scale" threshold. $3$ of $4$ frontier rows persist.}
\label{tab:frontier}
\footnotesize
\begin{tabular}{llccccc}
\toprule
Dataset & Model (scale) & faith$_{\textrm{base}}$ & faith$_{\textrm{v1}}$ & $\Delta_{\textrm{paradox}}$ & ref$_{7\textrm{B}}$ & persists? \\
\midrule
SQuAD     & Llama-3.3-70B  ($70$B)   & $0.767$ & $0.667$ & $\mathbf{+0.100}$ & $0.100$ & yes \\
SQuAD     & GPT-OSS-120B   ($120$B)  & $0.710$ & $0.680$ & $+0.030$ & $0.100$ & yes \\
PubMedQA  & Llama-3.3-70B  ($70$B)   & $0.552$ & $0.547$ & $+0.005$ & $0.025$ & no  \\
PubMedQA  & GPT-OSS-120B   ($120$B)  & $0.548$ & $0.536$ & $+0.013$ & $0.025$ & yes \\
\bottomrule
\end{tabular}
\end{table}

The headline finding: on SQuAD, Llama-$3.3$-$70$B exactly reproduces
the Mistral-$7$B paradox magnitude ($0.100$ vs $0.100$), and the
$120$B GPT-OSS model still exhibits a positive paradox of $0.030$.
HCPC-v$2$ recovery is also preserved at scale: $+0.091$ on SQuAD/$70$B
and $+0.026$ on SQuAD/$120$B, mirroring the $7$B recovery direction.

PubMedQA paradox magnitudes shrink at scale (from $0.025$ to
$0.005$/$0.013$). This is a prediction the coherence account makes:
stronger parametric priors are a more plausible substitute for
fragmented biomedical evidence than for short-span SQuAD answers,
where the missing connective evidence is the only path to the answer.
The qualitative pattern is preserved (HCPC-v$2$ still recovers
$+0.031$ faithfulness and \emph{halves} hallucination from $16.7\%$ to
$6.7\%$ on PubMedQA/$70$B), but the magnitude is attenuated.

This rules out the "small-model artifact" critique. The paradox is
not an emergent property of weak generators; it is a property of the
retrieval set's coherence that survives a $10\times$ scale increase in
generator parameters.

\subsection{Summary}
\label{sec:rob:summary}

Of the seven robustness checks: variance, RAPTOR (2/3 datasets), prompt
robustness, sub-chunk sensitivity (3/3 cells in the recovery
direction), and frontier-scale ($3/4$ rows persist; SQuAD/$70$B
exactly reproduces the $7$B magnitude) confirm the headline result.
The long-form check qualifies it as short-form-specific. The
noise-vs-coherence check rules out one form of the strongest version
of our claim (paradox magnitude exceeds noise) but is consistent with
a weaker, qualitative version (paradox is sign-distinct from noise on
the multi-hop and biomedical regimes). We incorporate both
qualifications in \S\ref{sec:limitations}.
